% short/06_results.tex: which model is best (Figure 8: one version of each of the six models of Section 5, mapped to
% their families in the caption; the CNN's lead stated through the paired comparisons, not through overlapping CIs) and
% the honest estimate; calibration before and after (Figure 9: gradient boosting, where it mattered, with the
% calibration slope that the eye can check, and the final CNN, whose raw outputs were already close to calibrated);
% several attempts (Figure 10); who it fails on (Figure 11); would more participants help (Figure 12, with how the
% extrapolation was computed: b in points and the curve's own value at 19, the base of the gain). Numbers only
% through \R keys.
\section{Results}\label{sec:results}

\begin{figure}[tbp]
\centering
\genfig{final/FINAL_15_families}
\caption{One version of each of the six models of Section~\ref{sec:methods} (\R{final/n_candidates} models were
compared in all), each tested on participants it had never seen (LOSO): the yaw baseline, the physics model, two
models on summary numbers (gradient boosting and the stack) and two that read the whole time series (MiniRocket and
the 1D CNN). Accuracy (left) and AUC (right) with their 95\,\% CIs. The 1D CNN, the final model (blue), scores
highest on both.}
\label{fig:families}
\end{figure}

\paragraph{Which model is best.} Every model separates the hands far better than chance, and every other model scores
higher than the yaw baseline (Figure~\ref{fig:families}). The yaw baseline itself does no better than the sign of the
first-half yaw alone (\R{describe/fhy_sign_rule_accuracy}, Section~\ref{sec:physics}): reading the yaw at more
moments adds little, and the models that do better also read other channels. The 1D CNN performed best: in its own
LOSO run (the slightly flattering score of Section~\ref{sec:methods}) its accuracy is \Rci{final/m3_deep_accuracy}
and its AUC \Rcib{final/m3_deep_auc}, against \Rci{final/ref_accuracy} and \Rcib{final/ref_auc} for the yaw baseline.
Its lead over the next models is small: compared with it on the same trials, participant by participant, the stack
and gradient boosting could not be told apart from it, while MiniRocket, the physics model and the yaw baseline did
detectably worse. Repeating the choice of the model inside every LOSO round gives the honest estimate for a new
participant: an accuracy of \Rci{final/nested_accuracy}, an AUC of \Rcib{final/nested_auc} and a log loss of
\Rcib{final/nested_log_loss}.

\begin{figure}[tbp]
\centering
\genfig{final/FINAL_16_calibration}
\caption{Calibration before and after, for a model that needed it (top: gradient boosting) and for the final model
(bottom: the 1D CNN). (a) Raw scores (the uncalibrated \pright{} $= \sigma(s)$) against the observed share of
right-hand trials, in \R{final/calfig_cnn_groups} groups of trials with similar scores (bars: 95\,\% CIs), with the
calibration slope; the dashed diagonal is perfect calibration. (b) The fitted calibration function, one curve per LOSO
round, each fitted on that round's training participants. (c) The calibrated probabilities against the observed
share, with the calibration slope.}
\label{fig:calibration}
\end{figure}

\paragraph{Calibration.} Gradient boosting, for example, was over-confident: groups of trials with an uncalibrated
\pright{} near 0 held more right-hand trials than claimed, and groups near 1 fewer (Figure~\ref{fig:calibration}, top left). One number
summarises this pattern, the \emph{calibration slope}: the $a$ of a single Platt function fitted to all the pooled
scores at once (1: nothing to correct; below 1: over-confident). For gradient boosting it is
\Rcib{final/calfig_gbm_raw_slope}. Its \R{final/calfig_gbm_folds} fitted functions, one per LOSO round ($a$ from
\R{final/calfig_gbm_a_min} to \R{final/calfig_gbm_a_max}), pull the scores towards 50\,\% (top middle). Afterwards the
groups scatter on both sides of the diagonal instead of lying above it at low scores and below it at high ones
(top right), and the slope is \R{final/calfig_gbm_platt_slope}; the log loss changes little
(\R{final/calfig_gbm_raw_ll} before, \R{final/calfig_gbm_platt_ll} after). The final CNN's raw outputs were already close to calibrated (bottom left; slope
\R{final/calfig_cnn_raw_slope}): its fitted functions stay near the diagonal ($a$ from \R{final/calfig_cnn_a_min} to
\R{final/calfig_cnn_a_max}), change its \pright{} by \R{final/calfig_cnn_change} percentage points on average and leave
its log loss practically unchanged (\R{final/calfig_cnn_raw_ll} before, \R{final/calfig_cnn_platt_ll} after). We kept
this step, planned for this model before the results, and the CNN's numbers in this report include it.

\begin{figure}[tbp]
\centering
\genfig{final/FINAL_17_attempts}
\caption{Combining $k$ attempts of the same participant and hand (the \R{final/agg_persons} participants with at
least \R{final/att_k_max} attempts of each hand; random draws of $k$ attempts, each participant left out in turn;
accuracy averaged over participants and hands): accuracy (left) and log loss (right, with 95\,\% bands), for the
combination of equation~\eqref{eq:aggregate} (blue, solid) and for a plain average of the probabilities (grey,
dashed).}
\label{fig:attempts}
\end{figure}

\paragraph{Several attempts.} Combining attempts with equation~\eqref{eq:aggregate} helps (Figure~\ref{fig:attempts}).
This test uses the \R{final/agg_persons} participants with at least \R{final/att_k_max} trials of each hand and
averages the accuracy over participants and hands, so a single attempt scores \R{final/att_k1_used_acc} here rather
than the \R{final/m3_deep_accuracy} above. The accuracy rises from \R{final/att_k1_used_acc} with one attempt to
\R{final/att_k3_used_acc} with three and \R{final/att_k5_used_acc} with five, and the log loss falls from
\R{final/att_k1_used_ll} to \R{final/att_k3_used_ll} and \R{final/att_k5_used_ll}. A plain average of the
probabilities of the attempts gives nearly the same answers, but its log loss levels off near
\R{final/att_k3_avg_ll} from three attempts on: an average does not become more confident as the evidence
accumulates. Like the CNN's own score above, these numbers are slightly flattering, because the CNN was chosen on the
same participants.

\begin{figure}[tbp]
\centering
\begin{minipage}[t]{0.49\linewidth}\centering
  \genfig{final/FINAL_18_persons}
  \caption{Accuracy of the final model for each participant when that participant was left out, sorted; dark bars:
  Android, hatched: iPhone; L: the two left-handed participants; dashed: \R{final/person_threshold}.}
  \label{fig:persons}
\end{minipage}\hfill
\begin{minipage}[t]{0.49\linewidth}\centering
  \genfig{final/FINAL_19_learning_curve}
  \caption{Accuracy of the 1D CNN trained on fewer participants (grey: one random draw; black: the mean of the
  draws), the fitted curve with its 95\,\% band, and its extrapolation to \R{final/lc_extrapolation_n}
  participants.}
  \label{fig:learning}
\end{minipage}
\end{figure}

\paragraph{Who it fails on.} The accuracy differs between participants (Figure~\ref{fig:persons}). The median
participant scores \R{final/brk_person_median}, \R{final/brk_n_perfect} participants have every trial right, and
\R{final/brk_below80} fall below \R{final/person_threshold}, the lowest being \R{final/brk_low_person_id} at
\R{final/brk_low_person} (\R{final/brk_low_person_n} trials). The two left-handed participants score
\R{final/brk_left_p11} (P11) and \R{final/brk_left_p17} (P17). The accuracy is \R{final/brk_android_acc} on Android
phones and \R{final/brk_ios_acc} on iPhones.

\paragraph{Would more participants help?} We retrained the CNN inside every LOSO round on random subsets of
\R{final/lc_sizes} of the training participants, drawing each size \R{final/lc_m3_deep_reps} times
(Figure~\ref{fig:learning}). The mean accuracy rises from \R{final/lc_m3_deep_acc_n03} with \R{final/lc_m3_deep_n0}
training participants to \R{final/lc_m3_deep_acc_n19} with all \R{final/lc_m3_deep_nall}, and flattens. We fitted a
curve that rises quickly and then levels off towards a ceiling $A$, the inverse power law
$\mathrm{accuracy}(n) = A - B\,n^{-C}$ for $n$ training participants ($A$ = \R{final/lc_m3_deep_asym},
$B$ = \Rb{final/lc_m3_deep_b_pts} percentage points, $C$ = \R{final/lc_m3_deep_c}), and read it at $n$ =
\R{final/lc_extrapolation_n}: \Rci{final/lc_m3_deep_at40}, only \Rb{final/lc_m3_deep_gain40}~points above the
curve's own value at \R{final/lc_m3_deep_nall} participants (\R{final/lc_m3_deep_fit_nall}). The interval, from
refitting the curve \R{final/lc_m3_deep_nboot} times on resampled draws, understates the uncertainty, because all
draws come from the same \R{describe/n_persons} people and the curve is extended beyond the measured sizes. More
participants of the same kind would therefore add little accuracy; new ones would matter mainly to check the accuracy
in groups we barely sampled, such as left-handed or older people and other phones, and for other kinds of movement.
